\section{Related Work}
\label{sec:related}

\textbf{Fraud detection with machine learning.}
Traditional approaches to insurance fraud detection usually rely on supervised classifiers such as logistic regression, decision trees, and ensemble methods \citep{phua2010comprehensive}. More recent work has looked at Bayesian belief networks for healthcare fraud \citep{kumaraswamy2024bayesian}, Bayesian methods for provider fraud detection \citep{ekin2013bayesian}, and reinforcement learning for real-time fraud detection \citep{saddi2024realtime}. Cost-sensitive learning has also been used to deal with the unequal costs of false positives and false negatives in fraud detection \citep{phua2010comprehensive}. However, these methods still make a decision on every case, without a principled way to express uncertainty or send unclear cases to human reviewers.

\textbf{Learning to defer and selective prediction.}
The idea of giving a classifier a reject option goes back to \citet{chow1957optimum}, and \citet{geifman2017selective} later extended it to deep networks with risk guarantees on the selected subset. \citet{madras2018predict} formalised learning to defer (L2D) as a broader version of rejection learning, where accuracy and fairness are optimised together. \citet{mozannar2020consistent} then established the Bayes-consistency of the $(K{+}1)$-class surrogate loss, which provides the theoretical basis for end-to-end L2D training. More recent extensions include probabilistic L2D with missing expert annotations and workload constraints \citep{nguyen2025probabilistic}, sequential L2D for temporal decision-making \citep{joshi2023sequential}, and active learning for more sample-efficient L2D \citep{charusaie2022sample}. \citet{mozannar2023exact} showed that exact L2D with halfspaces is NP-hard and proposed MILP relaxations with coverage constraints. The two-stage approach of training a classifier and then setting a rejection threshold is well-established in selective prediction \citep{geifman2017selective}; our specific contribution relative to this literature is the CVaR-penalised threshold selection, which targets tail risk rather than average error on retained predictions. We additionally combine Bayesian and discriminative uncertainty signals when making the deferral decision.

\textbf{Risk-sensitive machine learning.}
CVaR-based objectives have been used in areas such as portfolio optimisation \citep{rockafellar2000optimization}, risk-sensitive RL \citep{tamar2015optimizing}, and distributional dead-end detection \citep{killian2023risk}. These methods sit within the broader family of spectral risk measures \citep{acerbi2002spectral}, which assign weights to the loss quantile function through a distortion function. Even so, the connection between risk-sensitive learning and deferral has received relatively little attention. Our CVaR-guided threshold selection is meant to address that gap by directly penalising tail risk in the deferral objective.

\textbf{Off-policy evaluation.}
OPE methods, including importance sampling, direct methods, and doubly robust estimators \citep{dudik2011doubly}, are widely used to evaluate policies from observational data. We use OPE to assess the L2D fraud policy, drawing on techniques from the clinical RL literature \citep{tang2022leveraging}. At the same time, we note that our OPE demonstration is based on a simulated behaviour policy, which limits how directly the results carry over to real-world settings.